\begin{frame}{자주 묻는 질문 (1)}
  \textbf{Q. 경사하강법이 지역 최솟값에 갇힐 수 있나요?}
  \vskip0.4em
  \begin{itemize}
    \item 이 절의 $J(\boldsymbol{w})$ (평균제곱오차)는 $\boldsymbol{w}$ 에 대한
      \kb{볼록함수}라 \textbf{지역 최솟값이 곧 전역 최솟값} ---
      어디서 출발하든 결국 같은 최적점에 수렴.
    \item 지역 최솟값이 실제 골칫거리가 되는 건 \textbf{Week 9
      신경망처럼 비볼록 손실}을 쓸 때.
  \end{itemize}
  \vskip0.8em
  \textbf{Q. 매 스텝 데이터 전체를 다 봐야 하나요?}
  \vskip0.4em
  \begin{itemize}
    \item 위 코드 = 매 스텝 $m$ 개 전부로 그래디언트를 계산하는
      \kb{배치 경사하강법}(batch GD).
    \item 데이터가 수백만 개면 스텝 하나조차 느려지는데, 실전은
      작은 묶음(mini-batch)만 보는 \kb{확률적 경사하강법}(SGD)을
      훨씬 더 많이 쓴다 --- Week 9에서 변형들을 자세히.
  \end{itemize}
\end{frame}
